\section{Du saut spectral 4 aux symétries du Monstre : déploiement APP
\texorpdfstring{$2\to6\to54$}{2--6--54}, réseau de Leech,
\texorpdfstring{$V^\natural$}{V natural} et corridor Moonshine}
\label{sec:cierre-equivalencias-corredor-moonshine-20260825}

Le corridor n'est pas une succession de coïncidences numériques. Il part du bloc
central d'APP, conserve l'orientation somme--produit, se ramifie dans les
couplages APP--Fock et APP--Witt, et se coordonne de nouveau par la
graduation. L'ordre causal contraignant est
\[\begin{aligned}
 B_c=9P_++5P_-
 &\longrightarrow 4\longrightarrow2\longrightarrow6\longrightarrow54\\
 &\longrightarrow
 \left\{
 \begin{array}{l}
  \text{degré deux de Fock et }[q]t_3=54,\\
  \text{alignement APP--Witt et }N_\alpha\cong\Lambda_{24}
 \end{array}
 \right.\\
 &\longrightarrow V^\natural
 \longrightarrow
 \bigl(\operatorname{Aut},\operatorname{ch},g\mapsto T_g\bigr)
 \longrightarrow\text{Moonshine}.
\end{aligned}\]
Les deux branches intermédiaires possèdent des codomaines distincts. La clé commune \(54\)
est un invariant gradué transporté par le corridor ; ce n'est pas une application
linéaire d'un scalaire vers une algèbre d'opérateurs de vertex.

\subsection{Hiérarchie typée et \(2\) réalisations de l'indice \(54\)}

\begin{theorem}[Hiérarchie spectrale et déploiement APP]
\label{thm:cierre-vtm-jerarquia-app}
Soit \(B_c=7I+2R=9P_++5P_-\) le bloc central d'APP.\@ Si
\(M_\Pi,M_\Sigma\) sont les compressions modulaires définies, alors
\[
 9-5=4,
 \qquad (B_c)_{i\ne j}=2,
 \qquad \sum M_\Pi-\sum M_\Sigma=51-45=6,
\]
et le déploiement sur les neuf positions donne
\[
 9\cdot6=459-405=54.
\]
Les quatre entiers appartiennent respectivement au saut spectral, au
couplage local, à la compression modulaire et à l'intégration bidimensionnelle.
\end{theorem}

\begin{proof}
La décomposition spectrale de \(B_c\) donne \(9\) et \(5\), tandis que
\(7I+2R\) fixe le coefficient non diagonal. Les deux sommes agrégées sont \(51\)
et \(45\). Chaque entrée comprimée représente neuf positions, de sorte que le
déploiement multiplie la différence par neuf. Aucun des quatre entiers
ne s'obtient en identifiant rétrospectivement leurs types.
\end{proof}

La droite d'orientation \(o_{\Sigma\Pi}\), la parité de feuillet et la
normalisation triangulaire étant fixées, le couplage orienté
\[
 \mathcal H_{\Sigma,\Pi}^{(m)}:
 M_{\mathrm{bal}}^{C_3,-}
 \longrightarrow
 \operatorname{End}\!\bigl(\mathcal F_2(\mathcal V_m)\bigr)
 \otimes o_{\Sigma\Pi}
\]
satisfait
\[
 \operatorname{Tr}_{o,\mathcal F_2}
 \left(\mathcal H_{\Sigma,\Pi}^{(m)}(\nu)\Gamma_2(g_m)\right)
 =-\frac{3m\,c(\nu)}2.
\]
Pour \(m=12\) et \(c(\delta)=-3\), la trace vaut \(54\). D'autre part,
\[
 t_3(q)=q^{-1}\prod_{n\ge1}(1+q^n+q^{2n})^{-12}
       =q^{-1}-12+54q-76q^2-243q^3+\cdots,
\]
d'où \([q]t_3=54\).

\begin{theorem}[Cône de compatibilité de l'indice commun]
\label{thm:cierre-vtm-indice-comun-54}
Les structures de départ de chaque branche étant conservées, on a
\[
 \boxed{
 D_{\mathrm{APP}}
 =\operatorname{Tr}_{o,\mathcal F_2}
   \left(\mathcal H_{\Sigma,\Pi}^{(12)}(\delta)\Gamma_2(g_{12})\right)
 =[q]t_3
 =v_\alpha^2
 =\frac{196884}{5\cdot729+1}
 =54.}
\]
Cette égalité constitue un cône d'évaluations vers le même entier. Elle n'
identifie pas \(W_{\mathrm{APP}}\), l'espace de Fock, le réseau voisin et
\(V^\natural\).
\end{theorem}

\begin{proof}
Le recensement APP donne \(459-405=54\) ; la trace orientée donne \(54\) ; le développement
du produit cyclotomique donne \([q]t_3=54\) ; et le vecteur radial du voisin
satisfait \(v_\alpha^2=4^2\cdot2+11\cdot2=54\). Enfin,
\(5\cdot729+1=3646\) et \(54\cdot3646=196884\). Les égalités sont comparées
après la construction de chaque objet et ne constituent donc pas une sélection numérique du
corridor.
\end{proof}

\paragraph{Statut exact de la hiérarchie \(4\to2\to6\to54\) et de l'identification réticulaire.}
La hiérarchie \(4\to2\to6\to54\), la trace APP--Fock et l'extraction
\([q]t_3\) sont des identités internes exactes. L'identification du voisin à
Leech utilise la chambre, le vecteur de voisinage et la classification réticulaire
déclarés.
%
